% GENERATED FILE. Edit canonical lessons, then run npm run generate:books.
% canonical-source-hash: fnv1a64:5bb49d1ff183cab2
% canonical-lessons: ES-C661-recepcionista, ES-C661-botones, ES-C661-huesped, ES-C661-estancia, ES-C661-conserje

\chapter{En recepción}
\label{ch:es-en-recepcion}
\hlchaptermodality{661}

\begin{chapteropening}
\textbf{By the end of this chapter:} \emph{I can deal with hotel staff, from the receptionist to the concierge, and say how long I am staying.}

Complete the last lesson of chapter 661: I can deal with hotel staff, from the receptionist to the concierge, and say how long I am staying.
\end{chapteropening}

\section[el recepcionista]{el recepcionista — the receptionist}

\label{lesson:ES-C661-recepcionista}



\begin{quote}
Before the new one: say the Spanish for the guesthouse, then the Spanish for the hostel.
\end{quote}

\subsection*{You'll want to know: el recepcionista}
\textbf{el recepcionista} — \textquotedblleft{}the receptionist\textquotedblright{}.

\begin{quote}
El recepcionista me da la llave.
\end{quote}

\begin{quote}
\emph{The receptionist gives me the key.}
\end{quote}

One form for both: \emph{el recepcionista}, \emph{la recepcionista}.

\begin{cousinweb}
\emph{recepcionista} is built on \emph{recepción}, a learned borrowing from Latin \emph{receptiō} \textquotedblleft{}a receiving\textquotedblright{} (EtymDB), from \emph{recipere} \textquotedblleft{}to receive\textquotedblright{}.
\end{cousinweb}

\begin{grammarlens}[title={What you've built}]
One more word for telling, explaining and describing at B1.
\end{grammarlens}

\subsection*{Your turn}
\begin{itemize}
\raggedright
  \item \emph{Say it:} \emph{el recepcionista}
  \item \emph{Say it:} \emph{el recepcionista}, once more
  \item \emph{Say it:} a sentence of your own with \emph{el recepcionista}
  \item \emph{Recall:} say the Spanish for the guesthouse, then the Spanish for the hostel, then say \emph{el recepcionista} again
\end{itemize}

\begin{tcolorbox}[breakable,colback=teal!4,colframe=teal!35!black,title={Before you move on}]
What does el recepcionista mean? (\textquotedblleft{}The receptionist\textquotedblright{}.) Say it once more.
\end{tcolorbox}
\section[el botones]{el botones — the bellboy, the porter (in a hotel)}

\label{lesson:ES-C661-botones}



\begin{quote}
Before the new one: say the Spanish for the hostel, then the Spanish for the receptionist.
\end{quote}

\subsection*{You'll want to know: el botones}
\textbf{el botones} — \textquotedblleft{}the bellboy, the porter (in a hotel)\textquotedblright{}.

\begin{quote}
El botones sube las maletas.
\end{quote}

\begin{quote}
\emph{The bellboy takes the suitcases up.}
\end{quote}

One person, but the word ends in -s: \emph{el botones}, \emph{los botones}.

\begin{cousinweb}
\emph{botones} is built on \emph{botón} \textquotedblleft{}button\textquotedblright{}, which comes from French \emph{bouton}; the porter is named for the rows of buttons on his uniform.
\end{cousinweb}

\begin{grammarlens}[title={What you've built}]
One more word for telling, explaining and describing at B1.
\end{grammarlens}

\subsection*{Your turn}
\begin{itemize}
\raggedright
  \item \emph{Say it:} \emph{el botones}
  \item \emph{Say it:} \emph{el botones}, once more
  \item \emph{Say it:} a sentence of your own with \emph{el botones}
  \item \emph{Recall:} say the Spanish for the hostel, then the Spanish for the receptionist, then say \emph{el botones} again
\end{itemize}

\begin{tcolorbox}[breakable,colback=teal!4,colframe=teal!35!black,title={Before you move on}]
What does el botones mean? (\textquotedblleft{}The bellboy, the porter (in a hotel)\textquotedblright{}.) Say it once more.
\end{tcolorbox}
\section[el huésped]{el huésped — the guest (at a hotel or a house)}

\label{lesson:ES-C661-huesped}



\begin{quote}
Before the new one: say the Spanish for the receptionist, then the Spanish for the bellboy.
\end{quote}

\subsection*{You'll want to know: el huésped}
\textbf{el huésped} — \textquotedblleft{}the guest (at a hotel or a house)\textquotedblright{}.

\begin{quote}
Los huéspedes desayunan a las ocho.
\end{quote}

\begin{quote}
\emph{The guests have breakfast at eight.}
\end{quote}

\emph{Una casa de huéspedes} — a guesthouse.

\begin{cousinweb}
\emph{huésped} is inherited from Latin \emph{hospitem}, a form of \emph{hospes} \textquotedblleft{}guest; host\textquotedblright{} (EtymDB).
\end{cousinweb}

\begin{grammarlens}[title={What you've built}]
One more word for telling, explaining and describing at B1.
\end{grammarlens}

\subsection*{Your turn}
\begin{itemize}
\raggedright
  \item \emph{Say it:} \emph{el huésped}
  \item \emph{Say it:} \emph{el huésped}, once more
  \item \emph{Say it:} a sentence of your own with \emph{el huésped}
  \item \emph{Recall:} say the Spanish for the receptionist, then the Spanish for the bellboy, then say \emph{el huésped} again
\end{itemize}

\begin{tcolorbox}[breakable,colback=teal!4,colframe=teal!35!black,title={Before you move on}]
What does el huésped mean? (\textquotedblleft{}The guest (at a hotel or a house)\textquotedblright{}.) Say it once more.
\end{tcolorbox}
\section[la estancia]{la estancia — the stay; the room; the ranch (Southern Cone)}

\label{lesson:ES-C661-estancia}



\begin{quote}
Before the new one: say the Spanish for the bellboy, then the Spanish for the guest.
\end{quote}

\subsection*{You'll want to know: la estancia}
\textbf{la estancia} — \textquotedblleft{}the stay; the room; the ranch (Southern Cone)\textquotedblright{}.

\begin{quote}
Disfrute de su estancia.
\end{quote}

\begin{quote}
\emph{Enjoy your stay.}
\end{quote}

In Argentina and Uruguay \emph{una estancia} is a large farm.

\begin{cousinweb}
\emph{estancia} is built on \emph{estar} \textquotedblleft{}to be, to stay\textquotedblright{}, which is inherited from Latin \emph{stāre} \textquotedblleft{}to stand\textquotedblright{}.
\end{cousinweb}

\begin{grammarlens}[title={What you've built}]
One more word for telling, explaining and describing at B1.
\end{grammarlens}

\subsection*{Your turn}
\begin{itemize}
\raggedright
  \item \emph{Say it:} \emph{la estancia}
  \item \emph{Say it:} \emph{la estancia}, once more
  \item \emph{Say it:} a sentence of your own with \emph{la estancia}
  \item \emph{Recall:} say the Spanish for the bellboy, then the Spanish for the guest, then say \emph{la estancia} again
\end{itemize}

\begin{tcolorbox}[breakable,colback=teal!4,colframe=teal!35!black,title={Before you move on}]
What does la estancia mean? (\textquotedblleft{}The stay; the room; the ranch (Southern Cone)\textquotedblright{}.) Say it once more.
\end{tcolorbox}
\section[el conserje]{el conserje — the concierge, the porter; the caretaker}

\label{lesson:ES-C661-conserje}



\begin{quote}
Before the new one: say the Spanish for the guest, then the Spanish for the stay.
\end{quote}

\subsection*{You'll want to know: el conserje}
\textbf{el conserje} — \textquotedblleft{}the concierge, the porter; the caretaker\textquotedblright{}.

\begin{quote}
El conserje nos recomienda un restaurante.
\end{quote}

\begin{quote}
\emph{The concierge recommends a restaurant to us.}
\end{quote}

In a school or a block of flats, \emph{el conserje} is the caretaker.

\begin{cousinweb}
\emph{conserje} was borrowed from French \emph{concierge}. Where the French word comes from is not settled.
\end{cousinweb}

\begin{grammarlens}[title={What you've built}]
One more word for telling, explaining and describing at B1.
\end{grammarlens}

\subsection*{Your turn}
\begin{itemize}
\raggedright
  \item \emph{Say it:} \emph{el conserje}
  \item \emph{Say it:} \emph{el conserje}, once more
  \item \emph{Say it:} a sentence of your own with \emph{el conserje}
  \item \emph{Recall:} say the Spanish for the guest, then the Spanish for the stay, then say \emph{el conserje} again
\end{itemize}

\begin{tcolorbox}[breakable,colback=teal!4,colframe=teal!35!black,title={Before you move on}]
What does el conserje mean? (\textquotedblleft{}The concierge, the porter; the caretaker\textquotedblright{}.) Say it once more.
\end{tcolorbox}
